\section*{Bab \ref{ch:b1:systems-of-points} --- Sistem Titik; Pengantar Benda Tegar}

\begin{solution}{exo:b1:systems-of-points:1}
$x_G = (2.0 \times 0 + 3.0 \times 1.0)/5.0 = \qty{0.60}{m}$. Bumi--Bulan:
$384000/82 = \qty{4700}{km}$ dari pusat Bumi, di dalam Bumi
(\qty{6370}{km}): Bumi bergoyang terhadap sebuah titik \qty{1700}{km} di bawah
permukaannya.
\end{solution}

\begin{solution}{exo:b1:systems-of-points:2}
$4.0v = 0.010 \times 800$: $v = \qty{2.0}{m/s}$. Peluru $\tfrac12 \times 0.010
\times 800^2 = \qty{3.2}{kJ}$, senapan \qty{8}{J}: dari energi kimia
mesiunya; benda yang ringan mengambil hampir seluruhnya ($E_k = p^2/2m$
dengan $p$ yang sama).
\end{solution}

\begin{solution}{exo:b1:systems-of-points:3}
$v = 20000/2500 = \qty{8.0}{m/s}$; $E_k$: \qty{200}{kJ} sebelumnya, $\tfrac12 \times
2500 \times 64 = \qty{80}{kJ}$ sesudahnya: $60\%$ hilang.
\end{solution}

\begin{solution}{exo:b1:systems-of-points:4}
Ujung: $\int_0^L x^2(M/L)\dd x = ML^2/3$; pusat: $\int_{-L/2}^{L/2}x^2(M/L)
\dd x = ML^2/12$; Huygens dengan $d = L/2$: $ML^2/12 + ML^2/4 = ML^2/3$.
\end{solution}

\begin{solution}{exo:b1:systems-of-points:5}
$E_k = \tfrac12 Mv^2(1 + k)$: tabung $1.5$, gelang $2$ kali $\tfrac12 Mv^2$.
Pada laju yang sama gelangnya berenergi lebih besar: ia menggelinding lebih jauh di lantai datar
(dengan gesekan yang sama) dan memanjat lebih tinggi ($h = v^2(1 + k)/2g$).
\end{solution}

\begin{solution}{exo:b1:systems-of-points:6}
HCl: $\mu = 35m_H/36 = 0.972m_H$; atom: $\mu = m_e/(1 + 1/1836) =
0.99946m_e$. Frekuensinya $\propto 1/\sqrt\mu$: $+1.4\%$ bagi HCl, $+0.027\%$
bagi hidrogen (geseran isotop antara spektrum H dan D bertumpu padanya).
\end{solution}

\begin{solution}{exo:b1:systems-of-points:7}
$J = ML^2/3$, $d = L/2$: $T = 2\pi\sqrt{(ML^2/3)/(MgL/2)} = 2\pi\sqrt{2L/3g}
= \qty{1.64}{s}$; panjang setaranya $2L/3 = \qty{0.67}{m}$. Sebuah pukulan di
pusat perkusinya membuat batangnya berputar terhadap porosnya tanpa
reaksi impulsif di sana: pukullah bolanya di ``titik manis'' tongkatnya dan
tangannya tak merasa perih.
\end{solution}

\begin{solution}{exo:b1:systems-of-points:8}
$J = \tfrac12 MR^2 = \qty{1.0e-3}{kg.m^2}$; $C = 4\pi^2J/T^2 = \qty{9.9e-3}{N.m/rad}$.
\end{solution}

\begin{solution}{exo:b1:systems-of-points:9}
$\omega = \qty{314}{rad/s}$, $E = \tfrac12 J\omega^2 = \qty{4.9e5}{J}$;
$\Gamma = J\omega/t = \qty{314}{N.m}$; putarannya: $\tfrac12\omega t/2\pi = 250$.
\end{solution}

\begin{solution}{exo:b1:systems-of-points:10}
$a = g\sin\alpha/(1 + k)$: bola ($k = 2/5$) $0.714$, tabung ($1/2$)
$0.667$, gelang ($1$) $0.500$ kali $g\sin\alpha$: bolanya pertama, gelangnya terakhir;
massa dan jari-jari tak muncul di mana pun.
\end{solution}

\begin{solution}{exo:b1:systems-of-points:11}
Momentum: $0.010 \times 400 = 2.01\,v$: $v = \qty{2.0}{m/s}$; lalu energi:
$h = v^2/2g = \qty{0.20}{m}$. $E_k$: \qty{800}{J} sebelumnya, $\tfrac12 \times
2.01 \times 4.0 = \qty{4.0}{J}$ sesudahnya: $99.5\%$ hilang menjadi deformasi dan
kalor. Tumbukannya tak lenting --- energi tak kekal melewatinya,
momentumnya kekal (tegangan talinya tegak dan tumbukannya singkat).
\end{solution}

\begin{solution}{exo:b1:systems-of-points:12}
$m_1v_1 = m_1v_1' + m_2v_2'$ dan $m_1v_1^2 = m_1v_1'^2 + m_2v_2'^2$ memberikan
$v_1 + v_1' = v_2'$ (bagilah persamaan energinya dengan persamaan momentumnya),
lalu menyusullah rumus-rumusnya. Pecahan energi ke sasarannya $4m_1m_2/(m_1 + m_2)^2$:
proton $100\%$; karbon $4 \times 12/169 = 28\%$; timbal $4 \times 207/208^2 =
1.9\%$. Inti ringan mengambil energi neutronnya dalam beberapa tumbukan
(air, grafit); timbal akan menuntut ratusan.
\end{solution}

\begin{solution}{pb:b1:systems-of-points:1}
\textbf{1.} Pada sumbunya, di pusat cakramnya, menurut kesetangkupan.

\textbf{2.} Koenig: $E_k = \tfrac12 Mv_G^2 + E_k^*$, dan dalam \omterm{def:b1:systems-of-points:barycentric}{kerangka
barisentrumnya} yoyonya berputar terhadap sumbunya: $E_k^* = \tfrac12 J\omega^2$.

\textbf{3.} Cincin berjari-jari $\rho$, bermassa $2M\rho\,\dd\rho/R^2$: $J = \int_0^R
\rho^2\cdot 2M\rho\,\dd\rho/R^2 = MR^2/2 = 0.5 \times 0.050 \times 9\times10^{-4}
= \qty{2.25e-5}{kg.m^2}$.

\textbf{4.} Talinya tetap dan terurai dari porosnya: titik pada
porosnya yang bersentuhan dengannya sesaat diam, sehingga $v_G = r\omega$.

\textbf{5.} $E_k = \tfrac12 Mv_G^2 + \tfrac12(\tfrac12 MR^2)(v_G/r)^2 = \tfrac12
Mv_G^2(1 + R^2/2r^2)$; $R^2/2r^2 = 9/(2 \times 0.16) = 28$: faktor $29$ ---
hampir seluruh energinya ada pada putarannya.

\textbf{6.} Berat $Mg$ ke bawah, tegangan $T$ ke atas: $Ma_G = Mg - T$.

\textbf{7.} Hanya tegangannya yang bermomen terhadap sumbunya (\omterm{def:b1:angular-momentum:moment}{lengan gaya} $r$):
$J\dot\omega = Tr$.

\textbf{8.} $\dot\omega = a_G/r$: $T = Ja_G/r^2 = (R^2/2r^2)Ma_G$; lalu
$Ma_G = Mg - (R^2/2r^2)Ma_G$, $a_G = g/29 = \qty{0.34}{m/s^2}$.

\textbf{9.} $T = M(g - a_G) = 0.050 \times 9.47 = \qty{0.47}{N}$, $97\%$ dari
beratnya: talinya nyaris menahannya.

\textbf{10.} $t = \sqrt{2h/a_G} = \qty{2.4}{s}$; $v_G = a_Gt = \qty{0.82}{m/s}$.
Sebuah batu: \qty{0.45}{s} dan \qty{4.4}{m/s}.

\textbf{11.} $\omega = v_G/r = \qty{205}{rad/s} = \qty{1960}{rpm}$.

\textbf{12.} $\tfrac12 Mv_G^2 = \qty{0.017}{J}$ berbanding $\tfrac12 J\omega^2 =
\qty{0.47}{J}$: $28$ kali lebih banyak energi pada putarannya daripada pada jatuhnya.

\textbf{13.} $Mgh = \qty{0.49}{J}$; $\tfrac12 Mv_G^2 \times 29 = 0.5 \times 0.050
\times 0.67 \times 29 = \qty{0.49}{J}$. Tegangannya bekerja pada titik
talinya yang diam (talinya tak bergerak): dayanya nol.

\textbf{14.} Gaya penghentinya bekerja menyusuri talinya, melalui sumbunya:
tak bermomen terhadapnya; \omterm{def:b1:angular-momentum:L}{momentum sudutnya} $J\omega$ tak tersentuh. Hanya gaya
tangensial (gesekan pada porosnya) yang dapat memperlambat putarannya.

\textbf{15.} $T \approx Mv_G/\Delta t + Mg = 0.050 \times 0.82/0.005 + 0.49 =
\qty{8.7}{N}$, delapan belas kali beratnya --- dari sinilah sentakan yang dirasakan
jarinya.

\textbf{16.} $J\dot\omega = -\Gamma$: $t = J\omega/\Gamma = 2.25\times10^{-5} \times
205/10^{-4} = \qty{46}{s}$.

\textbf{17.} Seluruh energi putarannya $\tfrac12 J\omega^2 = \qty{0.47}{J}$ berubah
menjadi $Mgh'$: $h' = \qty{0.96}{m}$, hampir setinggi semula; dalam praktiknya kurang
dari itu karena gesekan porosnya selama ia tidur dan kerugian pada
sentakannya.

\textbf{18.} $a_G$ terarah ke bawah (yoyonya melambat dalam perjalanan naik) sementara
$v_G$ terarah ke atas: $T = M(g - a_G) < Mg$ --- lebih kecil daripada beratnya.

\textbf{19.} $r = R$: faktornya $1 + 1/2 = 1.5$, $a_G = 2g/3 = \qty{6.5}{m/s^2}$
--- jatuh cepat, sedikit berputar: poros yang tipislah yang membuat sebuah yoyo.

\textbf{20.} \omterm{ex:b1:systems-of-points:rolling}{Menggelinding tanpa tergelincir}: titik bahan di $C$ berkecepatan
nol; medan kecepatan sebuah \omterm{def:b1:systems-of-points:solid}{benda tegar} yang salah satu titiknya diam adalah
rotasi terhadap titik itu (sumbu melalui $C$, dengan $\omega$ yang sama).

\textbf{21.} $F$ bekerja di bagian bawah porosnya, setinggi $R - r$ di atas
mejanya, mendatar menuju tangannya: momennya $F(R - r)$ terhadap $C$ dengan
arah yang menggelindingkan gelendongnya menuju tangannya; berat dan
reaksinya melalui $C$ (tak bermomen).

\textbf{22.} $(J + MR^2)\dot\omega = F(R - r)$, $a_G = R\dot\omega = FR(R - r)/
(\tfrac32 MR^2) = 2F(R - r)/3MR > 0$: menuju tangannya.

\textbf{23.} Garis kerja $F$ melalui $C$ ketika talinya
dari bagian bawah porosnya bertemu mejanya di $C$: $\cos\beta = r/R$,
$\beta = 82^\circ$. Lebih curam: momennya terhadap $C$ membalik, gelendongnya menggelinding
menjauh (dan bagi $\beta$ yang cukup besar ia terangkat).

\textbf{24.} $a_G = 2 \times 0.20 \times 0.026/(3 \times 0.050 \times 0.030) =
\qty{2.3}{m/s^2}$; gesekannya dari $Ma_G = F - f$ (gesekannya melawan $F$):
$f = 0.20 - 0.050 \times 2.3 = \qty{0.085}{N}$; ia harus tinggal di bawah $\mu_sMg
\approx \qty{0.2}{N}$ (untuk $\mu_s = 0.4$), yakni $F \lesssim \qty{0.5}{N}$.

\textbf{25.} \omterm{def:b1:systems-of-points:cm}{Pusat massa}: translasi $G$ dan tegangannya;
\omterm{def:b1:angular-momentum:L}{momentum sudut} terhadap sebuah sumbu: putarannya dan arah gelindingnya;
energi: kelajuannya di bawah dan tinggi pemanjatannya.
\end{solution}
